\let\cleardoublepage\clearpage
\chapter{Химическая технология}
\ID{ҒТАМР 61.53.91}{}

\begin{header}
\swa{Е.~Аубакиров, Ж.~Ташмухамбетова, Ф.~Ахметова, С.~Сатаева, А.~Кенесова, А.~Темирова}{ПОЛИМЕР ҚАЛДЫҚТАРЫН ТЕРМОКАТАЛИТИКАЛЫҚ ӨҢДЕУ ПРОЦЕСІ ҮШІН ТАБИҒИ
ЦЕОЛИТ НЕГІЗІНДЕГІ КАТАЛИЗАТОРЛАРДЫҢ БЕЛСЕНДІЛІГІН ЗЕРТТЕУ}

\tsp{1}Е.~Аубакиров,
\tsp{1}Ж.~Ташмухамбетова,
\tsp{2}Ф.~Ахметова\envelope,
\tsp{2}С.~Сатаева,
\tsp{1}А.~Кенесова,
\tsp{3}А.~Темирова
\end{header}

\begin{affil}
\tsp{1}әл-Фараби атындағы Қазақ ұлттық университеті, Алматы, Қазақстан,

\tsp{2}Жәңгірхан атындағы Батыс Қазақстан аграрлық-техникалық университеті, Орал, Қазақстан,

\tsp{3} ЖШС«Munai Eco Engineering», Алматы, Қазақстан

\corrauthor{Корреспондент-автор: firuza.92@mail.ru.}
\end{affil}

Катализатор ретінде Тайжүзген кен орны табиғи цеолитін молибден
тұздарымен модификацияланған композитті материалдар пайдаланылды.
Оңтайлы реакция жағдайлары 450°C температурада және 0,5-0,6 МПа қысымда
анықталды. Сұйық фракцияларды талдау, құрамында 2,0\% Mo бар
катализатордың дистилляттың ең жоғары шығымын (61,55\% дейін) қамтамасыз
ететіні және алканды және ароматты көмірсутектердің түзілуіне ықпал
ететіндігі анықталды, нәтижесінде алынған өнімді одан әрі мотор отыны
мен химиялық шикізат өндірісінде пайдалану үшін перспективалы етеді.

Термокаталитикалық гидрлеу пластмасса және резеңке қалдықтарын қайта
өңдеудің экологиялық таза және технологиялық тұрғыдан мүмкін әдісі
ретінде өзінің тиімділігін дәлелдеді.

{\bfseries Түйін сөздер:} табиғи цеолит, полимер қалдықтары,
термокаталитикалық өңдеу, мазут, сканерлеуші микроскоп, көмірсутектер.

\begin{header}
STUDY OF THE ACTIVITY OF NATURAL ZEOLITE-BASED CATALYSTS FOR THERMOCATALITIC TREATMENT OF POLYMER WASTES

\tsp{1}Y.~Aubakirov,
\tsp{1}Zh.~Tashmukhambetova,
\tsp{2}F.~Akhmetova\envelope,
\tsp{2}S.~Satayeva,
\tsp{1}A.~Kenessova,
\tsp{3}A.~Temirova
\end{header}

\begin{affil}
\tsp{1}al-Farabi Kazakh National University, Almaty, Kazakhstan,

\tsp{2}Zhangirkhan West-Kazakhstan agrarian technical university, Uralsk, Kazakhstan,

\tsp{3}LLP «Munai Eco Engineering», Almaty, Kazakhstan,

\href{mailto:e-mailfiruza.92@mail.ru}{e-mail:firuza.92@mail.ru}
\end{affil}

This paper presents the results of a study of thermocatalytic
hydrogenation for the conversion of polymer waste into valuable liquid
hydrocarbon products. Composite materials modified with molybdenum salts
of natural zeolite from the Taizhuzgen deposit were used as catalysts.
Optimal reaction conditions were determined at a temperature of 450°C
and a pressure of 0.5-0.6 MPa. Analysis of liquid fractions revealed
that the catalyst containing 2.0\% Mo provides the highest distillate
yield (up to 61.55\%) and promotes the formation of alkane and aromatic
hydrocarbons, making the resulting product promising for further use in
the production of motor fuels and chemical raw materials.
Thermocatalytic hydrogenation has proven its effectiveness as an
environmentally friendly and technologically advanced method for
recycling plastic and rubber waste. Glass is one of the most important
materials widely used in the construction industry and everyday life.
The demand for Glass Products is growing every year, which, in turn,
leads the glass industry to increase production volumes and maintain
their quality. Scientific and technological advances in glass production
have played a huge role in expanding the scope of its application. New
production methods and approaches to improving the existing
technological processes appeared, and new areas of application of glass
began to open up. By changing the chemical composition of the product,
various types of glass were created. In this regard, it is important to
study the methods of obtaining glass with different properties. For the
production of various products, a study of its physico-chemical
properties was carried out, preparing a composition of colored glass in
laboratory conditions. According to the results of the study, it was
found that the resulting colored glass can be used in decorative
purposes.

{\bfseries Keywords:} natural zeolite, polymer waste, thermocatalytic
treatment, fuel oil, scanning microscope, hydrocarbons.

\begin{header}
ИССЛЕДОВАНИЕ АКТИВНОСТИ КАТАЛИЗАТОРОВ НА ОСНОВЕ ПРИРОДНЫХ ЦЕОЛИТОВ ДЛЯ ТЕРМОКАТАЛИТИЧЕСКОЙ ПЕРЕРАБОТКИ ПОЛИМЕРНЫХ ОТХОДОВ

\tsp{1}Е.~Аубакиров,
\tsp{1}Ж.~Ташмухамбетова,
\tsp{2}Ф.~Ахметова\envelope,
\tsp{2}С.~Сатаева,
\tsp{1}А.~Кенесова,
\tsp{3}А.~Темирова
\end{header}

\begin{affil}
\tsp{1}Казахский национальный университет им. аль-Фараби, Алматы, Казахстан,

\tsp{2}Западно-Казахстанский аграрно-технический университет им. Жангирхана, Уральск, Казахстан,

\tsp{3}ТОО «Munai Eco Engineering», г. Алматы, Казахстан,

e-mail: firuza.92@mail.ru
\end{affil}

В данной работе представлены результаты исследования процесса
термокаталитического гидрирования с целью переработки полимерных отходов
в ценные жидкие углеводородные продукты. В качестве катализаторов были
использованы композиционные материалы, модифицированные молибденовыми
солями природного цеолита месторождения Тайжузген. Определены
оптимальные условия реакции при температуре 450 °С и давлении 0,5-0,6
МПа. Анализ жидких фракций показал, что катализатор, содержащий 2,0\%
Мо, обеспечивает наибольший выход дистиллятов (до 61,55\%) и
способствует образованию алкановых и ароматических углеводородов, что
делает получаемый продукт перспективным для дальнейшего использования в
производстве моторных топлив и химического сырья. Термокаталитическое
гидрирование доказало свою эффективность как экологически чистый и
технологичный метод переработки отходов пластика и резины.

{\bfseries Ключевые слова:} природный цеолит, полимерные отходы,
термокаталитическая переработка, мазут, сканирующий микроскоп,
углеводороды.

\begin{multicols}{2}
{\bfseries Кіріспе.} Қазіргі заманның өзекті экологиялық проблемаларының
бірі - шиналар, пластмассалар, резеңке бұйымдары сияқты көмірсутектерге
бай қалдықтардың үздіксіз жиналуы. Бұл заттар табиғи ыдырау мен
биодеградацияға өте төзімді, бұл олардың қоршаған ортада ұзақ уақыт
сақталуына әкеледі. Демек, олар полигондарда және ашық жерлерде жиналып,
елеулі экологиялық қауіптер тудырады. Бұл проблеманы шешу тұрақты даму
үшін өте маңызды және қалдықтарды өңдеудің озық технологияларын жақсарту
бойынша мәселелерді қамтиды. Мұндай әдістер қалдықтарды қазандық
отындарын, көмірсутек мономерлерін, майлау материалдарын, шайырды және
коксты қоса алғанда, негізгі өнеркәсіптік және тұтыну тауарларын өндіру
үшін пайдаланылатын құнды қайталама шикізатқа айналдыруға мүмкіндік
береді {[}1-5{]}.

Осы саладағы зерттеулердің негізгі бағыты құрамында көміртегі бар
қалдықтарды басқару және қайта өңдеудің тиімді жүйелерін құру болып
табылады. Қазіргі уақытта шиналарды, резеңкелерді және пластикті қайта
өңдеу әдістерінің көпшілігі энергияны қалпына келтіруге бағытталған
төмен температуралық термиялық өңдеулерді немесе пиролизді пайдаланады.
Дегенмен, бұл әдістердің бірнеше кемшіліктері бар, өйткені нәтижесінде
пайда болатын жеңіл фракциялар мен қатты қалдықтар толық емес жану және
нашар отын сапасымен байланысты мәселелерді жиі көрсетеді. Бұл
кемшіліктерді жою үшін көмірсутектерді жоғары температурада
термокаталитикалық немесе термодеструктивті бөлу тиімдірек және ұтымды
балама ұсынады {[}6-10{]}. Бұл технологиялар сұйық көмірсутектер,
жанармай компоненттері, кокстелетін материалдар және шайыр тәрізді
қалдықтар сияқты химия өнеркәсібі үшін құнды шикізаттарды өндіруді
жеңілдетеді {[}11-13{]}.

Пластикалық қалдықтарды кәдеге жарату контекстінде бір сатылы
термокаталитикалық түрлендіру жоғары сапалы мотор отындарын өндіруде
ерекше тиімділікті көрсетеді {[}14{]}. Қалдықтарды өңдеудің әртүрлі
әдістерінің ішінде термокаталитикалық ыдырау перспективалы әдістердің
бірі ретінде ерекшеленеді, өйткені ол полимер қалдықтарын синтетикалық
отынға әрі қарай өңдеуге жарамды көмірсутекті фракцияларға айналдыруға
мүмкіндік береді {[}15-16{]}.

Зерттеудің мақсаты шиналар мен пластикалық материалдарды қайта өңдеу
үшін термокаталитикалық гидрлеу процестерін зерттеу және жақсарту болды.
Тәжірибелер шикізат ретінде полимер қалдықтары мен мазутты, катализатор
ретінде табиғи цеолит негізіндегі инновациялық композиттік
катализаторларды қолдану арқылы жүргізілді.

{\bfseries Материалдар мен әдістер.} Термокаталитикалық гидрлеу
реакцияларында шикізат ретінде «Bericap Kazakhstan» ЖШС шығарған сусын
бөтелкелерінің қақпақтарын өлшемдері 3-6 мм ұнтақтау арқылы дайындалған.
Барлық химиялық реагенттер қосымша тазартусыз пайдаланылды және барлық
ерітінділер дистилденген суда дайындалды.

Тайжүзген кен орнының табиғи цеолиті алдын ала ұнтақтаудан және өлшемі
0,25 мм молекулалық тор арқылы електен өткеннен кейін пайдаланылды.
Цеолит 18 сағат бойы су моншасында үздіксіз араластыру кезінде 0,1 М
NH₄Cl ерітіндісімен ион алмастыру әдісі арқылы өңдеуден өтті. Қышқылдық
модификациядан кейін цеолит сүзгіден өткізіліп, хлорид иондары толығымен
жойылғанша тазартылған сумен мұқият жуылды. Кейіннен NH₄⁺ иондарын H⁺
иондарымен алмастыру үшін 100-500 °C-қа дейінгі температурада ауа
ағынында 4-5 сағат бойы термиялық активтеу жүргізілді. Термоактивтеу
кезінде ауа ағынының жылдамдығы 500 мл/мин деңгейінде сақталды.

Алынған катализатордың бір бөлігі сутегі иондарын молибден және вольфрам
иондарымен ауыстыру үшін 80 °C температурада аммоний молибдаты мен
аммоний вольфрам ерітінділерімен қосымша сіңдірілген және
модификацияланған үлгілер 4 сағат бойы 100-600 °C температура
диапазонында күйдіруге ұшырады.

Пластикалық қоспаны 0,5 МПа қысыммен, 450°C дейінгі температурада және
15 минут бойы үздіксіз араластырумен периодты режимде жұмыс істейтін
қондырғыда жүргізілді. Қондырғы көлемі 150 мл болатын тот баспайтын
болаттан реактордан тұрды. Реактор трансформатор және амперметрмен
басқарылатын айнымалы ток қыздырғышы арқылы қыздырылды. Реактордағы
температураны бақылау хромель-тамшы термопара арқылы жүзеге асырылды.
Термокаталитикалық гидрлеу каталитикалық белсенділікті және
селективтілікті бағалау үшін орындалды. Процесс блогындағы қысым азотты
пайдаланып жасалды және манометрмен өлшенді. Полимердің (7 г),
катализатордың (0,28 г) және паста түзетін агенттің (7 г) қоспасы
реакторға салынды, содан кейін реактордың тығыздығы тексерілді, азотпен
тазартылды, 0,5-0,6 МПа қысым орнатылды, содан кейін қыздыру құрылғысы
қосылды.

Реактордың қыздыруы процесті бастағаннан кейін 30-40 минуттан кейін
(қысым төмендеген кезде) өшірілді және жүйе бөлме температурасына дейін
салқындатылды. Процесс барысында пайда болған газ қаныққан NaCl
ерітіндісі бар газометрде жиналды. Түзілген газ мөлшері қысым
айырмашылығынан есептелді.

Сұйық өнімді әртүрлі температуралық (0-180°C, 180-250°C, 250-320°C)
фракцияларға бөлу арқылы айдау қондырғыда жүргізілді. Ол үшін бөлме
температурасына дейін салқындатылатын реактордың форсункасына
қабылдағышы бар конденсатор қосылды. Айдау кезінде реактор қыздырылып,
әртүрлі температуралық фракциялар жиналды. Құрылғыны салқындағаннан
кейін қалған қалдық реактордан алынды. Әрбір тәжірибе бірдей жағдайларда
үш рет жүргізілді. Тәжірибе нәтижелері арасындағы ауытқу 2\%-дан аспады.

Әрбір фракцияның құрамы электронды соққы ионизациясымен сканерлеу
режимінде жұмыс істейтін Agilent 7890A/5975C (АҚШ) газ
хроматографиясы-масс-спектрометрі арқылы талданды.

{\bfseries Нәтижелер және талқылау.} Термокаталитикалық гидрлеу әдісімен
полимерлі қалдықтарды өңдеу үшін тиімді катализаторды анықтау сұйық
өнімдердің шығымдылығын талдау негізінде жүргізілді. «Тайжүзген» кен
орнының цеолиті негізінде дайындалған және құрамында 0,5--2,0\% молибден
(Мо) бар катализаторлардың қатысуымен жүргізілген термокаталитикалық
гидрлеу нәтижелері 1-кестеде келтірілген.

Жүргізілген тәжірибелер негізінде тиімді катализатор ретінде Тайжүзген
кен орны цеолиті негізіндегі 2,0\% Мо отырғызылған катализатор болып
табылды. Таңдау полимер қалдықтарын термокаталитикалық гидрлеу кезінде
алынған сұйық фракциялардың шығымы бойынша жүргізілді.

Термокаталитикалық гидрлеу арқылы алынған сұйық өнімдердің химиялық
құрамы мен концентрациясын анықтау үшін газ хроматографиялық талдау
жүргізілді. Масс-спектрометриялық әдісті қолдана отырып, қайнау
температурасы 0-180 °C, 180-250 °C және 250-320 °C болатын
фракциялардағы көмірсутектердің топтық құрамы анықталды.
\end{multicols}

\tcap{1-кесте. 0,5 - 2,0 \% Мо отырғызылған «Тайжүзген» кен орны
цеолиті қатысында жүргізілген полимерлі қалдықтарды термокатализдік
гидрогенизациялық өңдеу процестері}
\begin{longtblr}[
  label = none,
  entry = none,
]{
  width = \linewidth,
  colspec = {Q[170]Q[159]Q[112]Q[125]Q[125]Q[234]},
  cells = {c},
  cells = {font = \small},
  row{1} = {font=\bfseries\small},
  cell{2}{1} = {r=4}{},
  hlines,
  vlines,
}
Бастапқы шикізат           & Катализатор      & m(0-180°C), \% & m(180-250°C), \% & m(250-320°C), \% & Сұйық дистилляттардың шығымы, \% \\
Полимер қалдықтары + мазут & 0,5\% Mo/цеолит  & 13,21          & 14,11            & 7,27             & 32,89                            \\
                           & 1,0\% Mo+ цеолит & 17,52          & 8,79             & 6,58             & 34,59                            \\
                           & 1,5\% Mo+ цеолит & 15,07          & 21,44            & 6,74             & 53,25                            \\
                           & 2,0\% Mo+ цеолит & 12,44          & 19,12            & 30,0             & 61,55                            
\end{longtblr}

\tcap{2-кесте. 2\% Мо/цеолит қатысында полимер қалдықтарын
термокаталитикалық өңдеуден алынған сұйық фракциялардың топтық
көмірсутектік құрамы}
\begin{longtblr}[
  label = none,
  entry = none,
]{
  width = \linewidth,
  colspec = {Q[281]Q[188]Q[235]Q[235]},
  cells = {c},
  cells = {font = \small},
  row{1} = {font=\bfseries\small},
  row{2} = {font=\bfseries\small},
  cell{1}{1} = {r=2}{},
  cell{1}{2} = {c=3}{0.657\linewidth},
  hlines,
  vlines,
}
Көмірсутектердің топтық құрамы & 2,0 \%-дық Мо/цеолит катализаторы қатысында көмірсутектердің шығымы, \% &           &           \\
                               & 0-180°С                                                                 & 180-250°С & 250-320°С \\
Алкан                          & 33,18                                                                   & 37,45     & 54,91     \\
Изоалкан                       & 7,33                                                                    & 33,49     & 5,94      \\
Алкен                          & 6,07                                                                    & 4,45      & 3,28      \\
Циклоалкан                     & 22,01                                                                   & 12,73     & 6,63      \\
Циклоалкен                     & 5,35                                                                    & 1,43      & 0,42      \\
Ароматты көмірсутектер         & 22,07                                                                   & 6,18      & 27,13     
\end{longtblr}

\begin{multicols}{2}
Бастапқы және модифицирленген цеолиттің сапалық және сандық құрамы
рентгендік флуоресценциялық талдаудың көмегімен анықталды. Алынған
нәтижелер энергиядисперсиялық спектроскопия көмегімен алынған
нәтижелерге толығымен сәйкес келеді.

Бастапқы цеолитпен дайындалған үлгілердің морфологиясы және қышқылсыз
әдіспен активтелген цеолит электронды микроскоптың 10000 есе үлкейту
зерттелді. Көріп отырғанымыздай, активтелмеген және активтелген
цеолиттің (1-сурет) бетінде айқын кеуектілігі бар ұсақ дисперсті,
борпылдақ құрылымдар байқалады.
\end{multicols}

\begin{figs}[1-сурет. Активтелмеген (а) және активтелген (ә) цеолиттің микроэлектрондық фотосуреті]

\begin{minipage}{0.8\textwidth}
\begin{minipage}{0.48\linewidth}
\centering
\fig[0.95\linewidth]{c/image2}[а]
\end{minipage}
\hfill
\begin{minipage}{0.48\linewidth}
\centering
\begin{tblr}{
  width = \linewidth,
  colspec = {Q[369]Q[233]Q[233]},
  cells = {c},
  row{1} = {font=\bfseries},
  hlines, vlines,
}
Элемент & wt.\% & at.\% \\
C  & 4,91  & 8,61  \\
O  & 39,93 & 52,64 \\
Na & 1,51  & 1,39  \\
Al & 9,56  & 7,47  \\
Si & 30,98 & 23,26 \\
K  & 10,42 & 5,62  \\
Fe & 2,69  & 1,02  \\
\end{tblr}
\vspace{0.3cm}
\end{minipage}
\end{minipage}
\begin{minipage}{0.8\textwidth}
\begin{minipage}{0.48\linewidth}
\centering
\fig[0.95\linewidth]{c/image3}[ә]
\end{minipage}
\hfill
\begin{minipage}{0.48\linewidth}
\centering
\begin{tblr}{
  width = \linewidth,
  colspec = {Q[369]Q[233]Q[233]},
  cells = {c},
  row{1} = {font=\bfseries},
  hlines, vlines,
}
Элемент & wt.\% & at.\% \\
C  & 7,34  & 12,28 \\
O  & 43,41 & 54,53 \\
Mg & 0,79  & 0,65  \\
Al & 8,67  & 6,46  \\
Si & 31,44 & 22,49 \\
K  & 3,81  & 1,96  \\
Fe & 4,54  & 1,63  \\
\end{tblr}
\vspace{0.3cm}
\end{minipage}
\end{minipage}

\end{figs}

\begin{multicols}{2}
Бастапқы және активтелген цеолиттердің беткі морфологиясы 10000 есе
үлкейту бойынша сканерлеуші электрондық микроскопия көмегімен зерттелді.
Алынған микросуреттер активтелген цеолит құрылымында елеулі өзгерістерді
анықтайды. Активтелмеген цеолит пішіні дұрыс емес агрегатталған
бөлшектері бар салыстырмалы түрде тығыз, кеуекті құрылыммен сипатталады.
Сонымен қатар, активтелген үлгі неғұрлым дамыған кеуектілікті және
бетінде біркелкі таралған ұсақ дисперсті, борпылдақ бөлшектердің болуын
көрсетеді.

Байқалған құрылымдық өзгерістер NH₄Cl ерітіндісімен өңдеу бойынша,
сонымен қатар термоактивтеу кезінде кристалдық тордың ішінара
реконструкциясына байланысты болуы мүмкін. Микрокеуекті және мезокеуекті
құрылымның қалыптасуы катализатордың меншікті бетінің ұлғаюына ықпал
етеді және оның термокаталитикалық гидрлеу реакцияларындағы
белсенділігін арттырады. Осылайша, цеолиттің активтелуі беттің химиялық
құрамын өзгертіп қана қоймайды, сонымен қатар оның морфологиялық
сипаттамаларына айтарлықтай әсер етеді, активті металл орталықтарының
дисперсиясы және каталитикалық процестердің тиімді ағыны үшін оңтайлы
жағдайларды қамтамасыз етеді.

Полимер қалдықтарын термокаталитикалық өңдеу нәтижелері бойынша
катализаторды тереңірек талдау үшін 2\% Mo/цеолит катализаторы үшін
құрылымдық-морфологиялық ерекшеліктері зерттелді.2-суретте 2\%
Mo/цеолит катализаторына 1000x,5000x,10000х ұлғайтылған микроскоптың
микроэлектрондық фотосуреттері көрсетілген.

Микроскоптың ажыратымдылығын 1000х, 5000х және 10000х дейін ұлғайта
отырып, 2\% Мо/цеолит негізіндегі катализатордың беттік морфологиясының
біртектілігі анықталды, онда өлшемдері 1,36; 1,70 және 6,66 мкм-ден
752,0 және 979,1 нм-ге дейінгі бөлшектер бақыланады. Жалпы алғанда,
катализатор түйіршікті құрылым болып табылады, бөлшек өлшемдерінің өте
кең диапазоны бар. Активтелген цеолиттің беттік модификациясы кеуек
жиектерін тегістейді және кеуектер мөлшерінің гетерогенділігін
арттырады.
\end{multicols}

\begin{figs}[2-сурет.2\% Мо/цеолит катализаторының микроэлектрондық фотосуреттері: а) x1000; ә) x5000; б) x10000]
  \fig[0.32\textwidth]{c/image4}[а)]
  \fig[0.32\textwidth]{c/image5}[ә)]
  \fig[0.32\textwidth]{c/image6}[б)]
\end{figs}

\begin{multicols}{2}
Осылайша, цеолиттің молибден тұзымен модификациясы катализатордың
морфологиясына айтарлықтай әсер етуі мүмкін екендігі және осыған
байланысты алынған катализаторлар көмірсутектерді гидрогенизациялау
процесінде термокаталитикалық конверсиялау процесінде сұйық
фракциялардың шығымы мен құрамына әр түрлі әсер етуі мүмкін екендігі
анықталды.

{\bfseries Қорытынды.} Жүргізілген зерттеулер Тайжүзген кен орнының табиғи
цеолиті негізінде дайындалған және молибден тұздарымен модификацияланған
катализаторлар қатысында көмірсутектерге бай полимер қалдықтарын
термокаталитикалық гидрлеу процесінің жоғары тиімділігін дәлелдеді.
Процестің оңтайлы шарттары 450 °С температурада және 0,5-0,6 МПа қысымда
анықталды.

Эксперименттік деректерге сәйкес, құрамында 2,0 мас.\% Мо бар цеолит
катализаторы сұйық дистилляттардың ең жоғары шығымын-61,55\% қамтамасыз
ететіні анықталды, бұл 0,5-1,5 мас.\% Мо бар үлгілермен салыстырғанда
8-28\% жоғары көрсеткішке ие. Алынған сұйық өнімдердің
газ-хроматографиялық талдауы олардың құрамында алкандардың
(33,18-54,91\%) және ароматты көмірсутектердің (6,18-27,13\%) басым
екенін көрсетті, бұл өнімдерді мотор отындары мен химия өнеркәсібіне
арналған шикізат ретінде пайдалануға мүмкіндік береді.

Сканерлеуші электрондық микроскопия және элементтік талдау нәтижелері
цеолиттің NH₄Cl ерітіндісімен өңделуі мен молибденмен модификациялануы
нәтижесінде кеуектіліктің артуын, бөлшектердің дисперсиялануын және
катализатор бетінің морфологиялық біртектілігінің жақсарғанын көрсетті.
Бұл факторлар катализатордың белсенді орталықтарының санының артуына
және термокаталитикалық гидрлеу реакцияларындағы селективтіліктің
жоғарылауына ықпал етеді. Әзірленген технологияны пластикалық
қалдықтарды кәдеге жарату және екінші көмірсутек шикізатын өндірудің
экологиялық мәселесін шешудің перспективалық тәсілі ретінде ұсынуға
болады.

\emph{{\bfseries Қаржыландыру.} Бұл зерттеу Қазақстан Республикасы Ғылым
және жоғары білім министрлігінің Ғылым комитетінің қаржылық қолдауымен
жүзеге асырылды. (Грант №BR28713351 «Құрамында көміртегі бар қалдықтарды
энергия ресурстарына және сұранысқа ие өнімдерге кешенді өңдеу үшін
жоғары тиімді каталитикалық технологияны әзірлеу»).}
\end{multicols}

\begin{center}
{\bfseries Әдебиеттер}
\end{center}

\begin{refs}
1. Thew C.X.E., Lee Z.S., Srinophakun P., Ooi C.W. Recent advances and
challenges in sustainable management of plastic waste using
biodegradation approach//Bioresource Technology.-
2023.-Vol.374:128772.  DOI 10.1016/j.biortech.2023.128772.

2. Han D., Han W., Chen H. Pyrolysis of waste tires: A
review//Polymers.-2023.-Vol.15(7):1604. DOI 10.3390/polym 15071604.

3. Osman Y.Yanchen, Sharif H.Zein. Latest advances in waste plastic
pyrolytic catalysis //Processes. -2022.-Vol.10(4): 683. DOI
10.3390/pr10040683.

4. Afash H., Ozarisoy B., Altan H., Budayan C. Recycling of tire waste
using pyrolysis: An environmental
perspective//Sustainability.-2023.-Vol.15(19): 14178. DOI
10.3390/su151914178.

5. El-Wafai N.A., Farag M.I., Abdel-Basit H.M., et al. Eco-friendly
degradation of natural rubber powder waste using some microorganisms
with focus on antioxidant and antibacterial activities of biodegraded
rubber//Processes.-2023.-Vol.11(8):2350. DOI 10.3390/pr11082350.

6. Pal S., Kumar A.et al. Recent advances in catalytic pyrolysis of
municipal plastic waste for the production of Hydrocarbon
fuels//Processes.-2022.-Vol.10(8):1497. DOI 10.3390/pr10081497.

7. Achilias D.S., Ciprioti S.C. et al. Thermal and catalytic recycling
of plastics from waste electrical and electronic equipment-Challenges
and perspectives//Polymers.-2024.-Vol.16(17):2538. DOI
10.3390/polym16172538.

8. Biakhmetov B., et al. A review on catalytic pyrolysis of municipal
plastic waste//WIREs Energy Environ.-2023.Vol.12(6):e495. DOI
10.1002/wene.495.

9. Nkiawete M.M., et al. Thermo-Catalytic Decomposition Comparisons:
Carbon catalyst structure, hydrocarbon feed and
regeneration//Catalysts.-2023.-Vol.13(10):1382. DOI
10.3390/catal13101382.

10. An L., Kou Z., Li R., Zhao Z. Research Progress in Fuel Oil
Production by Catalytic Pyrolysis Technologies of Waste
Plastics//Catalysts.-2024.-Vol.14(3).P.212. DOI 10.3390/catal14030212.

11. Seliverstov E.S., Furda L.V., Lebedeva O.E. Thermocatalytic
Conversion of Plastics into Liquid Fuels over
Clays//Polymers.-2022.-Vol.14(10):2115. DOI 10.3390/polym14102115.

12. Cai W., Kumar R., Zheng Y., Zhu Z. et al. Exploring the Potential
of Clay Catalysts in Catalytic Pyrolysis of mixed plastic waste for
fuel and energy recovery//Heliyon.-2023.-Vol.9(12):e23140. DOI
10.1016/j.heliyon.2023.e23140.

13. Syamsiro M., Cheng S., Hu W., Saptoadi H. et al. Liquid and
Gaseous Fuel from Waste Plastics by Sequential Pyrolysis and Catalytic
Reforming Processes over Indonesian Natural Zeolite Catalysts//Waste
Technology.-2014.-Vol.2.(2)- P.44-51. DOI 10.14710/2.2.44-51.

14. Wei J., Liu J., Zeng W., Dong Z., et al. Catalytic hydroconversion
processes for upcycling plastic waste to fuels and
chemicals//Catalysis Science \&
Technology.-2023.-Vol.13(5).-P.1258-1280. DOI 10.1039/D2CY01886A.

15. Vural U.S., Yinanc A., SevindirH.C. Two-stage thermocatalytic
conversion of waste XLPE to diesel-like fuel//Serbian Chemical
Society.  -2024.-Vol.89(6). -P.921-937. DOI 10.2298/JSC230828035V.

16. Qian K., Tian W., Li W., Wu S., Chen D., Feng Y. Catalytic
Pyrolysis of Waste Plastics over Industrial Organic
Solid-Waste-Derived Activated Carbon: Impacts of Activation Agents//
Processes.-2022.-Vol.10.(12):2668. DOI 10.3390/pr10122668.
\end{refs}

\begin{info}
\hspace{1em}\emph{{\bfseries Авторлар туралы мәлімет}}

Аубакиров Е.А. - химия ғылымдарының докторы, профессор, әл-Фараби
атындағы Қазақ ұлттық университеті, Алматы, Қазақстан, е-mail:
miral.64@mail.ru;

Ташмухамбетова Ж.Х. - химия ғылымдарының кандидаты, доцент, әл-Фараби
атындағы Қазақ ұлттық университеті, Алматы, Қазақстан, е-mail:
zheneta@mail.ru;

Ахметова Ф.Ж. - PhD, аға оқытушы, Жәңгірхан атындағы Батыс Қазақстан
аграрлық-техникалық университеті, Орал, Қазақстан, е-mail:
firuza.92@mail.ru;

Сатаева С.С.- химия ғылымдарының кандидаты, қауымдастырылған профессор,
Жәңгірхан атындағы Батыс Қазақстан аграрлық-техникалық университеті,
Орал, Қазақстан, е-mail: sataeva\_safura@mail.ru;

Кенесова А. - оқытушы, әл-Фараби атындағы Қазақ ұлттық университеті,
Алматы, Қазақстан, е-mail: kenessova.aruzhan@gmail.com;

Темирова А. - PhD, «Munai Eco Engineering» ЖШС директоры, Алматы,
Қазақстан, е-mail: aika179@mail.ru.

\hspace{1em}\emph{{\bfseries Information about the authors}}

Aubakirov Ye.- doctor of Chemical Sciences,professor, al-Farabi Kazakh
national university, Almaty, Kazakhstan, е-mail:
miral.64@mail.ru;

Tashmukhambetova Zh.- candidate of Chemical Sciences, assoc. professor,
al-Farabi Kazakh national university, Almaty, Kazakhstan,
е-mail:zheneta@mail.ru;

Akhmetova F.- PhD, senior lecturer, Zhangirkhan West-Kazakhstan agrarian
technical university, Uralsk, Kazakhstan, e-mail:
firuza.92@mail.ru;

Satayeva S.S.- candidate of Chemical Sciences, ass.professor,
Zhangirkhan West-Kazakhstan agrarian technical university, Uralsk,
Kazakhstan, е-mail: sataeva\_safura@mail.ru;

Kenessova A. - lecturer, al-Farabi Kazakh national university, Almaty,
Kazakhstan, е-mail: kenessova.aruzhan@gmail.com;

Temirova A. - PhD, Director of LLP «Munai Eco Engineering», Almaty,
Kazakhstan, е-mail: aika179@mail.ru.
\end{info}
